\documentclass[12pt,a4paper]{article}

\usepackage[T2A]{fontenc}
\usepackage[utf8]{inputenc}
\usepackage[russian]{babel}
\usepackage{amsmath,amsthm,amssymb,mathtools}
\usepackage{geometry}
\usepackage{enumitem}
\usepackage{hyperref}
\usepackage{array}
\usepackage{longtable}

\geometry{margin=2.2cm}

\theoremstyle{definition}
\newtheorem{definition}{Определение}[section]
\newtheorem{remark}{Замечание}[section]
\newtheorem{example}{Пример}[section]

\theoremstyle{plain}
\newtheorem{theorem}{Теорема}[section]
\newtheorem{lemma}{Лемма}[section]
\newtheorem{corollary}{Следствие}[section]
\newtheorem{proposition}{Утверждение}[section]

\title{Билет №4 \\ \small Понятие сложности алгоритмов. Классы P и NP. Полиномиальная сводимость задач. Примеры NP-полных задач. Приближенные алгоритмы. Методы решения задач о выполнимости, об удовлетворении ограничений. Эволюционные алгоритмы. (ИТМО) \\ Понятие сложности алгоритмов. Классы P и NP. Полиномиальная сводимость задач. Теорема Кука об NP-полноте задачи выполнимости булевой формулы. Примеры NPполных задач, подходы к их решению. Субоптимальные решения неразрешимых проблем и вычислительно сложных задач. (СПбГУ) \\ Понятие сложности алгоритмов. Классы P, NP. Теорема Кука об NP-полноте задачи выполнимости булевой формулы. Примеры NP-полных задач. (СПбПУ)}
\author{Сериков Владислав Александрович}
\date{26 апреля 2026}


\begin{document}
\maketitle

\tableofcontents
\newpage

\section{Понятие сложности алгоритмов}

\subsection{Общая идея}
Теория сложности --- раздел, изучающий \textbf{разрешимые задачи}, но уже не с точки зрения существования алгоритма, а с точки зрения \textbf{затрат ресурсов} на решение.

Основные ресурсы:
\begin{itemize}
\item \textbf{время} --- сколько шагов делает вычислительная модель;
\item \textbf{память} --- сколько памяти используется.
\end{itemize}

\begin{definition}
Пусть $T:\mathbb N\to\mathbb N$ --- функция. Язык $L$ принадлежит классу
\[
DTIME(T(n)),
\]
если существует детерминированная машина Тьюринга, распознающая $L$ и работающая на входе длины $n$ не более чем за $cT(n)$ шагов для некоторой константы $c>0$.
\end{definition}

\begin{remark}
Здесь речь идёт о \emph{детерминированном} времени, поэтому используется обозначение $DTIME$.
\end{remark}

\subsection{Почему выделяют полиномиальное время}
Полиномиальное время считается границей \emph{практической вычислимости}. Если алгоритм имеет полиномиальную сложность, то его считают эффективным; если требуется экспоненциальное время, задача обычно рассматривается как практически трудная.

\section{Классы \texorpdfstring{$P$}{P}, \texorpdfstring{$NP$}{NP} и \texorpdfstring{$EXP$}{EXP}}

\subsection{Класс \texorpdfstring{$P$}{P}}
\begin{definition}
Класс $P$ определяется как
\[
P=\bigcup_{c\ge 1} DTIME(n^c).
\]
\end{definition}

То есть $P$ --- это множество языков, распознаваемых за \textbf{полиномиальное время} детерминированной машиной Тьюринга.

\subsection{Класс \texorpdfstring{$EXP$}{EXP}}
\begin{definition}
Класс $EXP$ определяется как
\[
EXP=\bigcup_{c\ge 1} DTIME(2^{n^c}).
\]
\end{definition}

Это класс задач, для решения которых достаточно \textbf{экспоненциального времени}.

\subsection{Класс \texorpdfstring{$NP$}{NP} через сертификаты}
\begin{definition}
Язык $L\subseteq\{0,1\}^*$ принадлежит классу $NP$, если существуют:
\begin{itemize}
\item многочлен $p(n)$;
\item работающая полиномиальное время машина Тьюринга $M$,
\end{itemize}
такие, что для любого слова $x$ длины $n$
\[
x\in L \iff \exists u\in\{0,1\}^{p(n)}:\ M(x,u)=1.
\]
\end{definition}

\begin{definition}
Слово $u$ называется \textbf{сертификатом} (или свидетельством). Машина $M$ называется \textbf{проверяющей машиной}.
\end{definition}

\subsection{Интуиция}
Класс $NP$ --- это задачи, для которых:
\begin{quote}
\emph{решение может быть проверено за полиномиальное время, если дан короткий сертификат.}
\end{quote}

Важно: не требуется уметь \emph{найти} сертификат быстро, нужно только уметь \emph{проверить} его быстро.

\subsection{Соотношения между классами}
\begin{theorem}
Имеют место вложения
\[
P\subseteq NP\subseteq EXP.
\]
\end{theorem}

\begin{proof}
\textbf{1) Докажем $P\subseteq NP$.}

Пусть $L\in P$. Тогда существует полиномиальная детерминированная машина Тьюринга $M$, распознающая $L$.

Возьмём в качестве сертификата пустое слово $\varepsilon$. Его длина равна нулю, а нулевой многочлен --- тоже многочлен. Тогда можно использовать ту же машину $M$ как проверяющую:
\[
x\in L \iff M(x,\varepsilon)=1.
\]
Значит, $L\in NP$.

\medskip
\noindent
\textbf{2) Докажем $NP\subseteq EXP$.}

Пусть $L\in NP$. Тогда для входа $x$ длины $n$ существует сертификат длины не более $p(n)$, где $p$ --- многочлен, и есть проверяющая машина $M$, работающая за полиномиальное время.

Можно перебрать \emph{все} двоичные слова длины $p(n)$:
\[
2^{p(n)}
\]
штук. Для каждого запускаем $M(x,u)$. Если хотя бы для одного $u$ получим $1$, то $x\in L$.

Число сертификатов экспоненциально, время проверки каждого --- полиномиально, значит общее время экспоненциально. Следовательно, $L\in EXP$.
\end{proof}

\subsection{Вопрос \texorpdfstring{$P=NP$}{P=NP}}
В настоящее время неизвестно, равны ли классы $P$ и $NP$.

\begin{remark}
Большинство специалистов считает, что
\[
P\ne NP,
\]
но доказательства этого факта пока нет.
\end{remark}

\subsection{Альтернативное определение \texorpdfstring{$NP$}{NP}}

\begin{definition}
Пусть $T:\mathbb N\to\mathbb N$. Язык $L$ принадлежит классу
\[
NTIME(T(n)),
\]
если существует недетерминированная машина Тьюринга, распознающая $L$ за время $O(T(n))$.
\end{definition}

\begin{theorem}
\[
NP=\bigcup_{c\ge 1} NTIME(n^c).
\]
\end{theorem}

\begin{proof}[Идея доказательства]
Недетерминированный выбор можно понимать как \emph{угадывание сертификата}, а дальнейшую работу --- как его проверку. Обратно, если есть сертификат и проверяющая машина, то недетерминированная машина может сначала угадать сертификат, а потом проверить его.
\end{proof}

\section{Полиномиальная сводимость задач}

\subsection{Определение сводимости по Карпу}
\begin{definition}
Язык $L$ \textbf{полиномиально сводится по Карпу} к языку $L'$ (обозначение $L\le_p L'$), если существует вычислимая за полиномиальное время функция
\[
f:\{0,1\}^*\to\{0,1\}^*
\]
такая, что для всех $x$
\[
x\in L \iff f(x)\in L'.
\]
\end{definition}

\subsection{Смысл определения}
Если $L\le_p L'$, то любую задачу о принадлежности слов языку $L$ можно решать через задачу для $L'$:
\begin{enumerate}
\item по входу $x$ вычислить $f(x)$;
\item решить задачу для $f(x)$ относительно $L'$;
\item выдать тот же ответ.
\end{enumerate}

\subsection{Корректность этой схемы}
\begin{proposition}
Если $L\le_p L'$ и язык $L'$ распознаётся за полиномиальное время, то язык $L$ тоже распознаётся за полиномиальное время.
\end{proposition}

\begin{proof}
По определению сводимости функция $f$ вычисляется за полиномиальное время. Затем алгоритм для $L'$ работает на слове $f(x)$ тоже за полиномиальное время. Композиция двух полиномиальных по времени процедур остаётся полиномиальной. Следовательно, задача для $L$ решается за полиномиальное время.
\end{proof}

\subsection{Транзитивность}
\begin{theorem}
Если
\[
L\le_p L' \quad \text{и}\quad L'\le_p L'',
\]
то
\[
L\le_p L''.
\]
\end{theorem}

\begin{proof}
Пусть $f$ --- сведение $L$ к $L'$, а $g$ --- сведение $L'$ к $L''$. Тогда композиция
\[
h(x)=g(f(x))
\]
вычисляется за полиномиальное время, так как композиция полиномов есть полином. Кроме того,
\[
x\in L \iff f(x)\in L' \iff g(f(x))\in L'',
\]
то есть
\[
x\in L \iff h(x)\in L''.
\]
Следовательно, $h$ является полиномиальным сведением.
\end{proof}

\section{NP-трудные и NP-полные задачи}

\begin{definition}
Язык $L'$ называется \textbf{NP-трудным}, если для любого языка $L\in NP$ выполнено
\[
L\le_p L'.
\]
\end{definition}

\begin{definition}
Язык $L'$ называется \textbf{NP-полным}, если:
\begin{enumerate}
\item $L'\in NP$;
\item $L'$ NP-труден.
\end{enumerate}
\end{definition}

\subsection{Смысл NP-полноты}
NP-полная задача --- это задача:
\begin{itemize}
\item сама принадлежащая $NP$;
\item настолько трудная, что к ней полиномиально сводится \emph{любая} задача из $NP$.
\end{itemize}

\subsection{Что если P=NP}
\begin{theorem}
Если хотя бы одна NP-полная задача принадлежит классу $P$, то
\[
P=NP.
\]
\end{theorem}

\begin{proof}
Пусть $L'$ --- NP-полная и при этом $L'\in P$. Тогда для любого $L\in NP$ имеем $L\le_p L'$. По доказанной выше корректности полиномиального сведения из разрешимости $L'$ за полиномиальное время следует разрешимость $L$ за полиномиальное время. Значит, всякий язык из $NP$ лежит в $P$, то есть $NP\subseteq P$. Так как всегда $P\subseteq NP$, получаем $P=NP$.
\end{proof}

\section{Примеры задач из \texorpdfstring{$NP$}{NP}}

Ниже перечислены задачи, которые в конспекте явно указаны как принадлежащие классу $NP$.

\subsection{\textsc{INDSET}}
\textbf{Постановка.} Дан граф $G$ и число $k$. Нужно определить, существует ли в $G$ независимое множество размера $k$.

\textbf{Сертификат.} Список из $k$ вершин.

\textbf{Проверка сертификата:}
\begin{enumerate}
\item проверить, что перечислены именно вершины графа;
\item проверить, что их ровно $k$;
\item для каждой пары вершин проверить отсутствие ребра между ними.
\end{enumerate}

\textbf{Минимальный пример.}
Пусть граф имеет вершины $\{1,2,3\}$ и единственное ребро $(1,2)$.
Тогда для $k=2$ сертификат $\{1,3\}$ корректен, потому что между $1$ и $3$ ребра нет.

\subsection{\textsc{TSP}}
\textbf{Постановка.} Дан взвешенный граф $G$ и число $k$. Нужно определить, существует ли цикл веса $k$, проходящий через все вершины.

\textbf{Сертификат.} Список вершин в порядке обхода цикла.

\textbf{Проверка сертификата:}
\begin{enumerate}
\item убедиться, что перечислены все вершины графа по одному разу;
\item проверить наличие нужных рёбер;
\item посчитать суммарный вес цикла;
\item сравнить его с $k$.
\end{enumerate}

\textbf{Минимальный пример.}
Треугольник с вершинами $1,2,3$ и весами всех рёбер, равными $1$. Для $k=3$ сертификат $(1,2,3)$ корректен.

\subsection{\textsc{GCONN}}
\textbf{Постановка.} Дан граф $G$ и вершины $u,v$. Нужно проверить, существует ли путь из $u$ в $v$.

\textbf{Сертификат.} Сам путь.

\textbf{Проверка сертификата:}
\begin{enumerate}
\item проверить, что первая вершина равна $u$, последняя --- $v$;
\item проверить, что каждая соседняя пара в списке соединена ребром.
\end{enumerate}

\textbf{Минимальный пример.}
Если в графе есть рёбра $(1,2)$ и $(2,3)$, то для $u=1$, $v=3$ сертификатом служит путь $(1,2,3)$.

\subsection{\textsc{GI}}
\textbf{Постановка.} Даны два графа $G_1$ и $G_2$. Нужно определить, изоморфны ли они.

\textbf{Сертификат.} Перестановка вершин.

\textbf{Проверка сертификата:}
\begin{enumerate}
\item применить перестановку к матрице смежности $G_1$;
\item сравнить полученную матрицу с матрицей смежности $G_2$.
\end{enumerate}

\subsection{\textsc{SUBSETSUM}}
\textbf{Постановка.} Дано множество целых чисел
\[
A=\{a_1,\dots,a_n\}
\]
и число $T$. Требуется определить, существует ли подмножество с суммой $T$.

\textbf{Сертификат.} Само подмножество.

\textbf{Проверка сертификата:}
\begin{enumerate}
\item проверить, что все элементы сертификата принадлежат $A$;
\item сложить их;
\item сравнить сумму с $T$.
\end{enumerate}

\textbf{Минимальный пример.}
Для $A=\{2,5,7\}$ и $T=9$ сертификат $\{2,7\}$ корректен.

\subsection{\textsc{COMPOSITE}}
\textbf{Постановка.} Дано целое число. Нужно определить, является ли оно составным.

\textbf{Сертификат.} Нетривиальный делитель.

\textbf{Проверка сертификата:}
\begin{enumerate}
\item убедиться, что делитель больше $1$ и меньше самого числа;
\item проверить деление без остатка.
\end{enumerate}

\textbf{Минимальный пример.}
Для числа $15$ сертификатом служит $3$, поскольку $15\bmod 3=0$.

\subsection{\textsc{FACTORING}}
\textbf{Постановка.} Даны числа $N,L,U$. Нужно определить, есть ли у $N$ простой делитель в промежутке от $L$ до $U$.

\textbf{Сертификат.} Такой простой делитель.

\subsection{\textsc{LP}}
\textbf{Постановка.} Дана система линейных неравенств с рациональными коэффициентами. Требуется определить, существует ли рациональное решение.

\textbf{Сертификат.} Подстановка значений переменных.

\textbf{Проверка сертификата:}
\begin{enumerate}
\item подставить значения в каждое неравенство;
\item проверить истинность всех неравенств.
\end{enumerate}

\subsection{\textsc{0/1-IP}}
\textbf{Постановка.} Та же задача, но неизвестные принимают только значения из $\{0,1\}$.

\textbf{Сертификат.} Двоичная подстановка.

\subsection{\textsc{SAT}}
\textbf{Постановка.} Дана булева формула. Нужно определить, выполнима ли она.

\textbf{Сертификат.} Подстановка значений переменных.

\textbf{Проверка сертификата:}
\begin{enumerate}
\item подставить значения в формулу;
\item вычислить значение формулы;
\item проверить, что получилась истина.
\end{enumerate}

\textbf{Минимальный пример.}
Для формулы
\[
(x\vee \neg y)\wedge y
\]
сертификат $x=1$, $y=1$ делает формулу истинной.

\section{Теорема Кука--Левина о NP-полноте \texorpdfstring{\textsc{SAT}}{SAT}}

\subsection{Формулировка}
\begin{theorem}[Кука--Левина]
Задача \textsc{SAT} является NP-полной.
\end{theorem}

Это центральный результат: \textsc{SAT} --- первая стандартная NP-полная задача.

\subsection{Структура доказательства}
Нужно доказать два пункта:
\begin{enumerate}
\item $\textsc{SAT}\in NP$;
\item любая задача из $NP$ полиномиально сводится к \textsc{SAT}.
\end{enumerate}

\subsection{Почему \textsc{SAT} принадлежит \texorpdfstring{$NP$}{NP}}
Сертификатом является подстановка значений переменных, а её проверка сводится к вычислению значения формулы за полиномиальное время.

\subsection{Идея трудной части доказательства}
Пусть $L\in NP$. Тогда существует:
\begin{itemize}
\item проверяющая машина $M$;
\item полиномиальный по длине входа сертификат $u$,
\end{itemize}
такие, что
\[
x\in L \iff \exists u\colon M(x,u)=1.
\]

Нужно по слову $x$ построить булеву формулу $\varphi_x$ такую, что
\[
x\in L \iff \varphi_x \text{ выполнима}.
\]

\subsection{Главная идея доказательства}
Доказываем сначала NP-полноту задачи \textsc{CSAT} (выполнимость формулы в КНФ), а затем получает NP-полноту \textsc{SAT} как следствие.

Ключевая идея:
\begin{quote}
нужно закодировать \textbf{корректное вычисление машины Тьюринга} булевой формулой.
\end{quote}

\subsection{Почему наивная идея не работает}
Можно было бы попытаться построить по $x$ булеву функцию
\[
f_x(u)=M(x,u),
\]
а затем записать её формулой в КНФ. Но эта функция зависит от всех сертификатов длины $p(n)$, их число равно $2^{p(n)}$, поэтому таблица истинности экспоненциальна.

\subsection{Локальность вычисления}
Правильная идея опирается на \textbf{локальность} работы машины Тьюринга:
\begin{itemize}
\item один шаг зависит только от текущего состояния;
\item символа(ов), которые сейчас читаются;
\item локального изменения на ленте.
\end{itemize}

Именно поэтому корректность вычисления можно описывать \emph{локальными ограничениями}.

\subsection{Шаги построения формулы}
Предлагаем считать, что машина:
\begin{itemize}
\item имеет отдельную входную полуленту;
\item отдельную рабочую полуленту;
\item является забывающей.
\end{itemize}

Далее вводится понятие \textbf{снэпшота} --- локального описания шага вычисления:
\[
[a,b,q],
\]
где:
\begin{itemize}
\item $a$ --- символ на входной ленте;
\item $b$ --- символ на рабочей ленте;
\item $q$ --- состояние.
\end{itemize}

\paragraph{Шаг 1.}
Для слова $x$ фиксированной длины $n$ рассматривается последовательность снэпшотов
\[
z_1,z_2,\dots,z_{T(n)},
\]
где $T(n)$ --- полиномиальная оценка времени работы машины.

\paragraph{Шаг 2.}
Вводятся булевы переменные, кодирующие:
\begin{itemize}
\item сам сертификат $y$;
\item все снэпшоты $z_1,\dots,z_{T(n)}$.
\end{itemize}

\paragraph{Шаг 3.}
Строится формула
\[
\varphi_x=(1)\wedge(2)\wedge(3)\wedge(4),
\]
где:
\begin{enumerate}
\item первые биты $y$ совпадают с входом $x$;
\item $z_1$ кодирует правильную начальную конфигурацию;
\item каждый следующий снэпшот корректно получается из предыдущего;
\item последний снэпшот является допускающим.
\end{enumerate}

\subsection{Почему эта формула корректна}
\begin{proposition}
Формула $\varphi_x$ выполнима тогда и только тогда, когда существует сертификат $u$, на котором проверяющая машина принимает $x$.
\end{proposition}

\begin{proof}
\textbf{Если} машина принимает $x$ с некоторым сертификатом $u$, то можно взять:
\begin{itemize}
\item код сертификата;
\item код всех реальных снэпшотов вычисления.
\end{itemize}
Тогда все четыре группы ограничений выполняются, значит $\varphi_x$ выполнима.

\medskip
\noindent
\textbf{Если} формула $\varphi_x$ выполнима, то значения её переменных задают:
\begin{itemize}
\item некоторый сертификат;
\item некоторую последовательность снэпшотов.
\end{itemize}
Условия (1)--(4) гарантируют, что это \emph{корректное принимающее вычисление} машины на входе $x$. Следовательно, существует сертификат, допускающий $x$, то есть $x\in L$.
\end{proof}

\subsection{Почему построение полиномиально}
Все ограничения локальны, число шагов вычисления полиномиально, размер описания одного шага постоянен относительно $n$. Поэтому:
\begin{itemize}
\item число переменных полиномиально;
\item число локальных ограничений полиномиально;
\item итоговая формула строится за полиномиальное время.
\end{itemize}

\subsection{Минимальная иллюстрация идеи}
Пусть машина за два шага проверяет очень короткое свойство. Тогда вместо общей таблицы истинности мы описываем:
\begin{itemize}
\item \emph{какой} была начальная конфигурация;
\item \emph{как} из неё получилась вторая;
\item \emph{почему} вторая является допускающей.
\end{itemize}
Это и есть принцип Кука--Левина: \textbf{вычисление превращается в систему булевых ограничений}.

\section{NP-полнота задачи \texorpdfstring{\textsc{3SAT}}{3SAT}}

\subsection{Определение}
\textsc{3SAT} --- это задача проверки выполнимости булевой формулы в \textbf{3-КНФ}, то есть в КНФ, где каждый дизъюнкт содержит ровно три различных литерала.

\begin{theorem}
Задача \textsc{3SAT} NP-полна.
\end{theorem}

\subsection{Схема доказательства}
Сводим к ней уже NP-полную задачу \textsc{CSAT}.

\subsection{Алгоритм преобразования дизъюнктов}
Пусть дан дизъюнкт. Нужно заменить его на эквивалентную по выполнимости формулу в 3-КНФ.

\paragraph{Случай 1: один литерал $x$.}
Заменяем на
\[
(x\vee u\vee v)\wedge(x\vee u\vee \neg v)\wedge(x\vee \neg u\vee v)\wedge(x\vee \neg u\vee \neg v).
\]

\textbf{Корректность.}
\begin{itemize}
\item Если $x=1$, то все четыре дизъюнкта истинны.
\item Если $x=0$, то при любых $u,v$ один из четырёх дизъюнктов будет ложным.
\end{itemize}

\textbf{Минимальный пример.}
Дизъюнкт $(x)$ превращается в четыре дизъюнкта выше.

\paragraph{Случай 2: два литерала $x\vee y$.}
Заменяем на
\[
(x\vee y\vee u)\wedge(x\vee y\vee \neg u).
\]

\textbf{Корректность.}
\begin{itemize}
\item Если $x\vee y=1$, оба дизъюнкта истинны.
\item Если $x=y=0$, то получаем $u\wedge\neg u$, что невыполнимо.
\end{itemize}

\textbf{Минимальный пример.}
\[
(x\vee y)\rightsquigarrow (x\vee y\vee u)\wedge(x\vee y\vee \neg u).
\]

\paragraph{Случай 3: три литерала.}
Ничего делать не нужно.

\paragraph{Случай 4: четыре литерала $x\vee y\vee z\vee u$.}
Заменяем на
\[
(x\vee y\vee w)\wedge(\neg w\vee z\vee u).
\]

\textbf{Корректность.}
\begin{itemize}
\item Если исходный дизъюнкт выполним, то один из двух новых можно удовлетворить за счёт уже истинного литерала, а второй --- правильным выбором $w$.
\item Если исходный дизъюнкт ложен, то все четыре литерала ложны, и никаким выбором $w$ оба новых дизъюнкта одновременно истинными сделать нельзя.
\end{itemize}

\paragraph{Общий случай: $k>4$.}
Дизъюнкт
\[
x_1\vee x_2\vee \dots \vee x_k
\]
разбивается рекурсивно:
\[
(x_1\vee \dots \vee x_{k-2}\vee w)\wedge(\neg w\vee x_{k-1}\vee x_k),
\]
после чего первый длинный дизъюнкт снова разбивается тем же способом.

\subsection{Почему преобразование полиномиально}
Каждый исходный дизъюнкт заменяется на формулу, размер которой растёт линейно по числу литералов. Следовательно, общее преобразование всей формулы выполняется за полиномиальное время.

\section{NP-полнота задачи \texorpdfstring{\textsc{0/1-IP}}{0/1-IP}}

\begin{theorem}
Задача \textsc{0/1-IP} NP-полна.
\end{theorem}

\subsection{Идея сведения}
Сводим \textsc{CSAT} к \textsc{0/1-IP}.

\subsection{Шаги преобразования}
Каждому булеву литералу ставим в соответствие линейное выражение:
\begin{itemize}
\item положительный литерал $u$ заменяем на переменную $u$;
\item отрицательный литерал $\neg u$ заменяем на $1-u$.
\end{itemize}

Если дизъюнкт имеет вид
\[
u_1\vee \neg u_3\vee \neg u_{15},
\]
то он преобразуется в неравенство
\[
u_1+(1-u_3)+(1-u_{15})\ge 1.
\]

\subsection{Корректность}
\begin{proposition}
Дизъюнкт выполним тогда и только тогда, когда соответствующее неравенство имеет решение в $\{0,1\}$.
\end{proposition}

\begin{proof}
Сумма в левой части равна числу истинных литералов данного дизъюнкта. Неравенство
\[
\text{сумма}\ge 1
\]
выполняется тогда и только тогда, когда истинным является хотя бы один литерал, то есть тогда и только тогда, когда дизъюнкт истинен.
\end{proof}

\subsection{Минимальный пример}
\[
(x\vee \neg y)\rightsquigarrow x+(1-y)\ge 1.
\]
Проверим:
\begin{itemize}
\item если $x=0$, $y=1$, то $0+(1-1)=0$, неравенство ложно;
\item если $x=1$, $y=1$, то $1+0=1$, верно;
\item если $x=0$, $y=0$, то $0+1=1$, верно.
\end{itemize}
Это exactly соответствует истинности дизъюнкта.

\section{NP-полнота задачи \texorpdfstring{\textsc{INDSET}}{INDSET}}

\begin{theorem}
Задача \textsc{INDSET} NP-полна.
\end{theorem}

\subsection{Сведение из \textsc{3SAT}}
Пусть дана формула в 3-КНФ с $m$ дизъюнктами.

\subsection{Шаги построения графа}
\begin{enumerate}
\item Для каждого дизъюнкта создаём \textbf{кластер} из $7$ вершин.
\item Каждая вершина кластера соответствует одной из $7$ подстановок трёх переменных, удовлетворяющих дизъюнкту.
\item Внутри кластера все вершины соединяем попарно.
\item Между вершинами разных кластеров проводим ребро, если соответствующие частичные подстановки \textbf{несогласованы}, то есть назначают одной и той же переменной разные значения.
\end{enumerate}

\subsection{Почему именно 7 вершин}
У дизъюнкта из трёх литералов всего $2^3=8$ подстановок, и ровно одна из них делает дизъюнкт ложным. Значит, удовлетворяющих подстановок $7$.

\subsection{Корректность сведения}
\begin{proposition}
Формула выполнима тогда и только тогда, когда в построенном графе существует независимое множество размера $m$.
\end{proposition}

\begin{proof}
\textbf{Если формула выполнима.}
Берём общую удовлетворяющую подстановку и для каждого дизъюнкта выбираем вершину, соответствующую сужению этой подстановки на переменные данного дизъюнкта. Получаем по одной вершине из каждого кластера. Они попарно согласованы, поэтому рёбер между ними нет. Значит, это независимое множество размера $m$.

\medskip
\noindent
\textbf{Если есть независимое множество размера $m$.}
Так как внутри каждого кластера все вершины соединены, можно взять не более одной вершины из кластера. Чтобы размер был $m$, нужно взять ровно по одной вершине из каждого кластера. Отсутствие рёбер между выбранными вершинами означает согласованность соответствующих частичных подстановок. Значит, их можно объединить в общую подстановку, которая удовлетворяет каждому дизъюнкту и всей формуле.
\end{proof}

\subsection{Минимальная иллюстрация}
Если формула состоит из двух дизъюнктов, то строятся два кластера по $7$ вершин. Независимое множество размера $2$ должно взять по одной вершине из каждого кластера, причём эти две вершины не должны конфликтовать по значениям общих переменных.

\section{NP-полнота задачи \texorpdfstring{\textsc{TSP}}{TSP}}

\begin{theorem}
Задача \textsc{TSP} NP-полна.
\end{theorem}

Конспект доказывает это не напрямую, а через цепочку сведений:
\[
\textsc{CSAT}\le_p \textsc{dHAMPATH}\le_p \textsc{HAMPATH}\le_p \textsc{HAMCYCLE}\le_p \textsc{TSP}.
\]

\subsection{\textsc{CSAT} \texorpdfstring{$\le_p$}{<=p} \textsc{dHAMPATH}}
По формуле в КНФ строится ориентированный граф. Основная идея:
\begin{itemize}
\item каждой переменной соответствует цепочка, которую можно пройти слева направо или справа налево;
\item направление обхода кодирует истинность переменной;
\item вершины, соответствующие дизъюнктам, вставляются как ``экскурсии'' из этих цепочек;
\item экскурсия возможна именно тогда, когда соответствующий литерал удовлетворяет дизъюнкту.
\end{itemize}

\subsubsection*{Шаги построения}
\begin{enumerate}
\item Для каждой переменной строится длинная ориентированная цепочка.
\item Добавляются стартовая и конечная вершины.
\item Для каждого дизъюнкта добавляется отдельная вершина.
\item Если переменная входит в дизъюнкт положительно, делается экскурсия, согласованная с обходом слева направо.
\item Если отрицательно --- экскурсия строится в противоположном направлении.
\end{enumerate}

\subsubsection*{Корректность}
\begin{itemize}
\item Если формула выполнима, цепочки проходятся в соответствии со значениями переменных, а в каждый дизъюнкт можно сделать экскурсию через удовлетворяющий литерал.
\item Если в графе есть гамильтонов путь, направление обхода каждой цепочки задаёт значения переменных, а наличие посещения всех вершин-дизъюнктов гарантирует удовлетворимость всех дизъюнктов.
\end{itemize}

\subsection{\textsc{dHAMPATH} \texorpdfstring{$\le_p$}{<=p} \textsc{HAMPATH}}
Каждая вершина ориентированного графа заменяется тройкой вершин
\[
v^{(0)},\ v^{(1)},\ v^{(2)},
\]
соединённых в цепочку. Направленная дуга $u\to v$ кодируется ребром между $u^{(2)}$ и $v^{(0)}$.

\subsubsection*{Корректность}
Гамильтонов путь в ориентированном графе эквивалентен гамильтонову пути в неориентированном построенном графе: структура троек заставляет проходить вершины в правильном порядке, а переходы между тройками кодируют ориентированные дуги.

\subsection{\textsc{HAMPATH} \texorpdfstring{$\le_p$}{<=p} \textsc{HAMCYCLE}}
К графу добавляется новая вершина, соединённая со всеми исходными вершинами.

\subsubsection*{Корректность}
\begin{itemize}
\item Гамильтонов путь в исходном графе замыкается через новую вершину в гамильтонов цикл.
\item Гамильтонов цикл в новом графе, проходящий через новую вершину, после удаления этой вершины превращается в гамильтонов путь исходного графа.
\end{itemize}

\subsection{\textsc{HAMCYCLE} \texorpdfstring{$\le_p$}{<=p} \textsc{TSP}}
Всем рёбрам исходного графа приписывается вес $1$. Тогда вопрос:
\begin{quote}
``есть ли гамильтонов цикл?''
\end{quote}
равносилен вопросу:
\begin{quote}
``есть ли цикл веса $k$, проходящий через все вершины?''
\end{quote}
где $k$ --- число вершин.

\subsubsection*{Минимальный пример}
Если граф --- треугольник, то гамильтонов цикл существует, и в \textsc{TSP} это соответствует циклу веса $3$.

\section{Соотношение задач распознавания и поиска}

В конспекте отдельно отмечено, что для ряда NP-задач из умения отвечать ``да/нет'' можно получить и поиск решения.

\begin{theorem}
Если
\[
P=NP,
\]
то существует полиномиальный алгоритм, который по выполнимой булевой формуле находит удовлетворяющую подстановку.
\end{theorem}

\subsection{Алгоритм}
Пусть дана выполнимая формула $\varphi(x_1,\dots,x_n)$.

\begin{enumerate}
\item Для переменной $x_1$ проверяем, выполнима ли формула $\varphi[x_1:=0]$.
\item Если выполнима, фиксируем $x_1=0$, иначе $x_1=1$.
\item Переходим к $x_2$ и аналогично проверяем выполнимость уже частично подставленной формулы.
\item Продолжаем до $x_n$.
\end{enumerate}

\subsection{Корректность}
На каждом шаге хотя бы одно из двух значений сохраняет выполнимость формулы, иначе исходная формула была бы невыполнима. Поэтому процедура последовательно строит полную удовлетворяющую подстановку.

\subsection{Минимальный пример}
Пусть
\[
\varphi(x,y)=x\vee y.
\]
\begin{itemize}
\item Проверяем $\varphi[x:=0]=y$. Она выполнима.
\item Фиксируем $x=0$.
\item Проверяем оставшуюся формулу $y$:
значит $y=1$.
\end{itemize}
Получаем решение $(x,y)=(0,1)$.

\section{Класс \texorpdfstring{$coNP$}{coNP}}

\begin{definition}
\[
coNP=\{L:\ \overline L\in NP\}.
\]
\end{definition}

\subsection{Примеры}
\begin{itemize}
\item \textsc{UNSAT} --- множество невыполнимых формул;
\item \textsc{TAUTO} --- множество тавтологий.
\end{itemize}

\begin{theorem}
Задачи \textsc{SAT} и \textsc{TAUTO} являются coNP-полными.
\end{theorem}

\begin{remark}
Этот материал не является центральным для формулировок билетов, но естественно продолжает тему классов сложности и присутствует в конспекте.
\end{remark}

\section{Пространственная сложность: кратко}

Так как после NP-тем вводятся ещё и пространственные классы, кратко перечислим главное.

\begin{definition}
$PSPACE$ --- класс языков, распознаваемых детерминированной машиной Тьюринга с полиномиальным объёмом памяти.

$NPSPACE$ --- аналогично для недетерминированной машины.
\end{definition}

\begin{theorem}[Сэвитча]
\[
PSPACE=NPSPACE.
\]
\end{theorem}

\begin{theorem}
\[
P\subseteq NP,\ coNP\subseteq PSPACE\subseteq EXP\subseteq NEXP.
\]
\end{theorem}

\section{Подходы к решению вычислительно сложных задач}

\subsection{1. Сведение одной задачи к другой}
Это основной инструмент теории сложности.
Если задача $A$ сводится к задаче $B$, то:
\begin{itemize}
\item полиномиальный алгоритм для $B$ дал бы полиномиальный алгоритм для $A$;
\item неразрешимость/трудность $A$ передаётся задаче $B$.
\end{itemize}

\subsection{2. Проверка решения вместо его поиска}
Это и есть логика класса $NP$: решение трудно найти, но можно быстро проверить.

\subsection{3. Переход от задачи поиска к задаче разрешения}
Для некоторых задач поиска, в частности для \textsc{SAT}, поиск сводится к повторным вызовам решателя задачи разрешения.

\subsection{4. Экспоненциальный перебор сертификатов}
Из доказательства $NP\subseteq EXP$ следует базовый универсальный подход:
\begin{enumerate}
\item перебрать все сертификаты;
\item проверить каждый полиномиально.
\end{enumerate}
Этот подход всегда корректен для NP-задач, но обычно неэффективен.

\section{Подробные алгоритмы и их минимальные примеры}

Ниже собраны алгоритмы, которые прямо присутствуют в материале конспекта или немедленно следуют из него.

\subsection{Алгоритм проверки сертификата для \textsc{SAT}}
\textbf{Вход:} булева формула $\varphi$, подстановка $\sigma$.

\textbf{Шаги:}
\begin{enumerate}
\item Для каждой переменной формулы взять значение из $\sigma$.
\item Подставить значения в литералы.
\item Вычислить значения всех подформул снизу вверх.
\item Получить значение всей формулы.
\item Если значение равно $1$, принять сертификат; иначе отвергнуть.
\end{enumerate}

\textbf{Корректность.}
По определению выполнимости булева формула выполнима на данной подстановке тогда и только тогда, когда после подстановки её значение равно истине. Алгоритм вычисляет именно это значение.

\textbf{Пример.}
\[
\varphi=(x\vee \neg y)\wedge y,\qquad \sigma(x)=1,\ \sigma(y)=1.
\]
Тогда:
\[
x\vee\neg y=1\vee 0=1,\qquad 1\wedge 1=1.
\]
Сертификат принимается.

\subsection{Алгоритм сведения \textsc{CSAT} к \textsc{3SAT}}
\textbf{Вход:} формула в КНФ.

\textbf{Шаги:}
\begin{enumerate}
\item Рассмотреть каждый дизъюнкт по отдельности.
\item Если в нём 1 литерал --- заменить по схеме из 4 дизъюнктов.
\item Если 2 литерала --- заменить по схеме из 2 дизъюнктов.
\item Если 3 литерала --- оставить как есть.
\item Если больше 3 литералов --- рекурсивно разбивать, вводя новые переменные.
\item Объединить результаты конъюнкцией.
\end{enumerate}

\textbf{Корректность.}
Для каждого типа дизъюнкта выполнимость сохраняется; значит сохраняется и выполнимость всей формулы.

\textbf{Пример.}
\[
(x\vee y)\rightsquigarrow (x\vee y\vee u)\wedge(x\vee y\vee \neg u).
\]

\subsection{Алгоритм сведения \textsc{CSAT} к \textsc{0/1-IP}}
\textbf{Вход:} формула в КНФ.

\textbf{Шаги:}
\begin{enumerate}
\item Для каждой булевой переменной ввести переменную $u_i\in\{0,1\}$.
\item Каждый положительный литерал заменить на $u_i$.
\item Каждый отрицательный литерал заменить на $1-u_i$.
\item Для каждого дизъюнкта записать неравенство ``сумма литералов $\ge 1$''.
\item Получить систему линейных неравенств.
\end{enumerate}

\textbf{Корректность.}
Каждый дизъюнкт истинен тогда и только тогда, когда хотя бы один его литерал равен $1$, то есть сумма не меньше $1$.

\textbf{Пример.}
\[
(x\vee \neg y\vee z)\rightsquigarrow x+(1-y)+z\ge 1.
\]

\subsection{Алгоритм построения графа для сведения \textsc{3SAT} к \textsc{INDSET}}
\textbf{Вход:} формула в 3-КНФ с $m$ дизъюнктами.

\textbf{Шаги:}
\begin{enumerate}
\item Для каждого дизъюнкта построить кластер из 7 вершин.
\item Каждой вершине сопоставить удовлетворяющую дизъюнкт частичную подстановку.
\item Внутри каждого кластера соединить все вершины попарно.
\item Между кластерами соединять только конфликтующие подстановки.
\item В качестве параметра для \textsc{INDSET} взять число $m$.
\end{enumerate}

\textbf{Корректность.}
Доказана выше: независимое множество размера $m$ эквивалентно выбору по одной согласованной удовлетворяющей подстановке на каждый дизъюнкт.

\subsection{Алгоритм получения удовлетворяющей подстановки из решателя \textsc{SAT}}
\textbf{Предположение:} $P=NP$, значит есть полиномиальный решатель \textsc{SAT}.

\textbf{Шаги:}
\begin{enumerate}
\item Для первой переменной попробовать значение $0$ и спросить, выполнима ли формула после такой подстановки.
\item Если да --- зафиксировать $0$, иначе зафиксировать $1$.
\item Повторить для следующей переменной.
\item После $n$ шагов получить полную подстановку.
\end{enumerate}

\textbf{Корректность.}
На каждом шаге хотя бы одно значение должно сохранять выполнимость, иначе исходная формула была бы невыполнима.

\textbf{Пример.}
\[
\varphi=(x\vee y)\wedge(\neg x\vee z).
\]
\begin{itemize}
\item Проверяем $\varphi[x:=0]=y\wedge z$; если выполнима, берём $x=0$.
\item Затем определяем $y$ и $z$ аналогично.
\end{itemize}

\section{Приближённые алгоритмы}

\subsection{Мотивация}
Для многих NP-трудных задач неизвестны полиномиальные алгоритмы точного решения, и, более того, широко распространено убеждение, что таких алгоритмов не существует, если $P\ne NP$. Поэтому естественно ослабить требование к качеству ответа и искать не точное, а \textbf{приближённое} решение, которое:
\begin{itemize}
\item вычисляется за полиномиальное время;
\item гарантированно не слишком далеко от оптимума.
\end{itemize}

Такой подход особенно важен для задач оптимизации:
\begin{itemize}
\item задачи коммивояжёра;
\item упаковки;
\item покрытия;
\item расписаний;
\item размещения объектов;
\item маршрутизации.
\end{itemize}

\subsection{Оптимизационные задачи и мера качества}
Пусть дана задача оптимизации. Для каждого допустимого решения $S$ определена целевая функция $f(S)$.

Различают:
\begin{itemize}
\item \textbf{задачи минимизации}: нужно минимизировать $f(S)$;
\item \textbf{задачи максимизации}: нужно максимизировать $f(S)$.
\end{itemize}

Обозначим:
\[
OPT(I)
\]
--- значение оптимального решения для экземпляра $I$,

\[
A(I)
\]
--- значение решения, найденного алгоритмом $A$.

\begin{definition}
Для задачи минимизации алгоритм называется \textbf{$\rho$-приближённым}, если для любого входа $I$
\[
A(I)\le \rho\cdot OPT(I), \qquad \rho\ge 1.
\]
\end{definition}

\begin{definition}
Для задачи максимизации алгоритм называется \textbf{$\rho$-приближённым}, если для любого входа $I$
\[
A(I)\ge \frac{1}{\rho}\,OPT(I), \qquad \rho\ge 1.
\]
\end{definition}

Эквивалентно обе ситуации можно записать единой оценкой:
\[
\max\left\{\frac{A(I)}{OPT(I)},\frac{OPT(I)}{A(I)}\right\}\le \rho.
\]

\subsection{Основные классы приближимости}

\begin{definition}
Алгоритм называется \textbf{приближённым с постоянным коэффициентом}, если $\rho$ есть константа, не зависящая от размера входа.
\end{definition}

\begin{definition}
Схема приближения \textbf{PTAS} (Polynomial Time Approximation Scheme) --- это семейство алгоритмов, которое для любого $\varepsilon>0$ строит решение с коэффициентом
\[
1+\varepsilon
\]
для задач минимизации (или $(1-\varepsilon)$ для максимизации), причём при фиксированном $\varepsilon$ время работы полиномиально по размеру входа.
\end{definition}

\begin{definition}
Схема приближения \textbf{FPTAS} (Fully Polynomial Time Approximation Scheme) --- это PTAS, у которой время работы полиномиально не только по размеру входа, но и по $1/\varepsilon$.
\end{definition}

\subsection{Общая схема доказательства оценки}
Типичная структура доказательства корректности приближённого алгоритма такова:
\begin{enumerate}
\item доказать, что алгоритм всегда строит допустимое решение;
\item оценить значение решения алгоритма сверху или снизу;
\item оценить оптимум через ту же величину;
\item сравнить две оценки и получить коэффициент приближения.
\end{enumerate}

\subsection{Жадные приближённые алгоритмы}
Один из самых распространённых подходов --- \textbf{жадный выбор}: на каждом шаге алгоритм принимает локально лучшее решение.

Плюсы:
\begin{itemize}
\item простота;
\item высокая скорость;
\item часто хорошие практические результаты.
\end{itemize}

Минусы:
\begin{itemize}
\item жадность редко даёт оптимум для NP-трудных задач;
\item для оценки качества обычно требуется отдельное доказательство.
\end{itemize}

\subsection{Пример: задача вершинного покрытия}
Напомним задачу \textsc{Vertex Cover}: по графу $G=(V,E)$ требуется найти минимальное множество вершин $C\subseteq V$, такое что каждое ребро инцидентно хотя бы одной вершине из $C$.

\subsubsection{Алгоритм}
\textbf{Вход:} граф $G=(V,E)$.

\textbf{Шаги:}
\begin{enumerate}
\item Инициализировать $C:=\varnothing$.
\item Пока есть непокрытое ребро $(u,v)$:
\begin{enumerate}
\item добавить в $C$ обе вершины $u$ и $v$;
\item удалить из рассмотрения все рёбра, инцидентные $u$ или $v$.
\end{enumerate}
\item Вернуть $C$.
\end{enumerate}

\subsubsection{Минимальный пример}
Пусть граф --- путь из трёх вершин:
\[
1-2-3.
\]
Если выбрать ребро $(1,2)$, алгоритм добавит вершины $1,2$. После этого оба ребра уже покрыты. Получаем
\[
C=\{1,2\}.
\]
Оптимальное решение:
\[
C^*=\{2\}.
\]
Здесь приближение равно $2$.

\subsubsection{Доказательство корректности}
На каждом шаге выбирается ребро $(u,v)$, и обе его вершины добавляются в покрытие. Значит, это ребро точно оказывается покрытым. После удаления всех рёбер, инцидентных $u$ и $v$, инвариант сохраняется: все удалённые рёбра покрыты. Когда алгоритм завершается, непокрытых рёбер не остаётся, значит, $C$ --- корректное вершинное покрытие.

\subsubsection{Оценка качества}
\begin{theorem}
Описанный алгоритм является $2$-приближённым для задачи \textsc{Vertex Cover}.
\end{theorem}

\begin{proof}
Рассмотрим множество рёбер, выбранных алгоритмом на разных шагах. Они попарно не имеют общих концов, то есть образуют паросочетание $M$.

Любое вершинное покрытие должно содержать хотя бы одну вершину из каждого ребра паросочетания, следовательно,
\[
OPT\ge |M|.
\]
Алгоритм добавляет по две вершины на каждое ребро из $M$, значит,
\[
|C|=2|M|.
\]
Следовательно,
\[
|C|=2|M|\le 2OPT.
\]
Значит, коэффициент приближения равен $2$.
\end{proof}

\subsection{Пример: задача покрытия множества}
Пусть дано универсальное множество
\[
U=\{e_1,\dots,e_n\}
\]
и семейство подмножеств
\[
\mathcal S=\{S_1,\dots,S_m\}, \qquad S_i\subseteq U.
\]
Нужно выбрать минимальное число множеств из $\mathcal S$, покрывающих все элементы $U$.

\subsubsection{Жадный алгоритм}
На каждом шаге выбирается множество, покрывающее наибольшее число ещё не покрытых элементов.

\subsubsection{Шаги}
\begin{enumerate}
\item $C:=\varnothing$, $R:=U$ --- множество ещё не покрытых элементов.
\item Пока $R\ne\varnothing$:
\begin{enumerate}
\item выбрать $S_i$, максимизирующее $|S_i\cap R|$;
\item добавить $S_i$ в решение $C$;
\item заменить $R:=R\setminus S_i$.
\end{enumerate}
\item Вернуть $C$.
\end{enumerate}

\subsubsection{Минимальный пример}
Пусть
\[
U=\{1,2,3,4\},
\]
\[
S_1=\{1,2\},\quad S_2=\{2,3,4\},\quad S_3=\{4\}.
\]
Тогда жадный алгоритм сначала выберет $S_2$, так как оно покрывает три элемента, потом $S_1$ для покрытия элемента $1$. Получится покрытие из двух множеств:
\[
\{S_2,S_1\}.
\]

\subsubsection{Замечание}
Для этой задачи жадный алгоритм не обязательно оптимален, но имеет логарифмическую оценку качества.

\subsection{Другие типовые техники приближения}
\begin{enumerate}
\item \textbf{Локальный поиск.}
Строится некоторое допустимое решение, затем выполняются локальные улучшения, пока они возможны.

\item \textbf{Округление решений линейных релаксаций.}
Сначала задача решается как линейная программа без целочисленности, затем дробное решение округляется.

\item \textbf{Динамическое программирование с масштабированием.}
Применяется в задачах типа \textsc{Knapsack}; позволяет получать FPTAS.

\item \textbf{Разделяй и властвуй + склейка решений.}

\item \textbf{Рандомизация.}
Случайный выбор может упростить алгоритм и улучшить ожидаемое качество.
\end{enumerate}

\section{Методы решения задач удовлетворения ограничений (CSP)}

\subsection{Общая постановка CSP}
\begin{definition}
\textbf{Задача удовлетворения ограничений} (Constraint Satisfaction Problem, CSP) задаётся:
\begin{itemize}
\item набором переменных
\[
X=\{x_1,\dots,x_n\};
\]
\item для каждой переменной $x_i$ --- доменом допустимых значений $D_i$;
\item набором ограничений $C=\{C_1,\dots,C_m\}$, каждое из которых запрещает или разрешает некоторые совместные значения нескольких переменных.
\end{itemize}
Нужно найти такое присваивание значений всем переменным, которое удовлетворяет всем ограничениям.
\end{definition}

\subsection{Примеры CSP}
\begin{enumerate}
\item \textbf{Булева выполнимость.}
Переменные булевы, ограничения --- дизъюнкты формулы.

\item \textbf{Раскраска графа.}
Переменные --- вершины, домены --- цвета, ограничения --- соседние вершины должны иметь разные цвета.

\item \textbf{Судоку.}
Переменные --- клетки, домены --- цифры от $1$ до $9$, ограничения --- различие в строках, столбцах и блоках.

\item \textbf{Расписание.}
Переменные --- занятия, домены --- временные слоты, ограничения --- отсутствие конфликтов.
\end{enumerate}

\subsection{CSP как общая модель}
CSP удобно тем, что многие задачи можно описать единым языком:
\begin{itemize}
\item переменные;
\item допустимые значения;
\item ограничения на совместимость.
\end{itemize}
После этого можно применять универсальные методы решения.

\subsection{Полный перебор}
Самый прямой метод --- перебрать все возможные присваивания.

\subsubsection{Алгоритм}
\begin{enumerate}
\item Выписать все комбинации
\[
(x_1,\dots,x_n)\in D_1\times\dots\times D_n.
\]
\item Для каждой комбинации проверить все ограничения.
\item При первой подходящей комбинации остановиться.
\end{enumerate}

\subsubsection{Корректность}
Алгоритм перебирает все возможные полные присваивания, значит, если решение существует, оно будет найдено. Если решение не найдено, значит, допустимого присваивания нет.

\subsubsection{Недостаток}
Если домены велики, полный перебор имеет экспоненциальную сложность.

\subsection{Поиск с возвратом (backtracking)}
Это базовый общий метод решения CSP.

\subsubsection{Идея}
Переменные присваиваются по одной. Как только частичное присваивание нарушает хотя бы одно ограничение, ветвь поиска отсекается.

\subsubsection{Алгоритм}
\textbf{Вход:} CSP с переменными $x_1,\dots,x_n$.

\textbf{Шаги:}
\begin{enumerate}
\item Если все переменные уже назначены, вернуть найденное решение.
\item Выбрать следующую переменную $x_i$.
\item Для каждого значения $a\in D_i$:
\begin{enumerate}
\item временно положить $x_i:=a$;
\item проверить, не нарушены ли ограничения на уже назначенных переменных;
\item если не нарушены, рекурсивно продолжить поиск;
\item если рекурсивный вызов нашёл решение, вернуть его.
\end{enumerate}
\item Если ни одно значение не подошло, выполнить откат и вернуть неуспех.
\end{enumerate}

\subsubsection{Минимальный пример}
Рассмотрим CSP:
\[
x\in\{1,2\},\qquad y\in\{1,2\},
\]
ограничение:
\[
x\ne y.
\]

Поиск:
\begin{itemize}
\item пробуем $x=1$;
\item пробуем $y=1$ --- нарушение;
\item пробуем $y=2$ --- подходит.
\end{itemize}
Получаем решение:
\[
x=1,\ y=2.
\]

\subsubsection{Доказательство корректности}
Рассуждаем по индукции по глубине рекурсии.
\begin{itemize}
\item На каждом шаге алгоритм перебирает все возможные значения выбранной переменной.
\item Отсечение происходит только тогда, когда уже нарушено какое-либо ограничение. Такое частичное присваивание не может быть продолжено до решения.
\item Значит, ни одно потенциально корректное полное решение не теряется.
\item Следовательно, если решение существует, одна из ветвей рекурсии приведёт к нему.
\item Если все ветви исчерпаны без успеха, решения нет.
\end{itemize}

\subsection{Проверка согласованности частичного присваивания}
Ключевой подалгоритм в backtracking --- быстрая проверка, не противоречит ли текущее частичное присваивание ограничениям.

\textbf{Шаги:}
\begin{enumerate}
\item Рассмотреть все ограничения, затрагивающие только уже назначенные переменные.
\item Проверить, выполняются ли они.
\item Если хотя бы одно нарушено, вернуть ``ложь''.
\item Иначе вернуть ``истина''.
\end{enumerate}

\subsection{Forward checking}
Это усиление backtracking.

\subsubsection{Идея}
После присваивания переменной значения нужно не только проверить уже нарушенные ограничения, но и удалить из доменов соседних переменных значения, которые стали невозможны.

\subsubsection{Шаги}
\begin{enumerate}
\item Назначить переменной значение.
\item Для каждой ещё не назначенной переменной, связанной ограничением:
\begin{enumerate}
\item удалить из её домена несовместимые значения;
\item если домен опустел, немедленно сделать откат.
\end{enumerate}
\end{enumerate}

\subsubsection{Преимущество}
Противоречие обнаруживается раньше, чем в обычном backtracking.

\subsection{Поддержание согласованности дуг}
Для бинарных CSP часто используют \textbf{согласованность дуг}.

\begin{definition}
Ограничение между переменными $x$ и $y$ согласовано по дуге $x\to y$, если для каждого значения $a\in D_x$ существует хотя бы одно значение $b\in D_y$, совместимое с $a$.
\end{definition}

\subsubsection{Идея метода}
Если значение $a\in D_x$ не имеет ни одного допустимого ``партнёра'' в $D_y$, то его можно удалить из домена $x$.

\subsubsection{Общий алгоритм}
\begin{enumerate}
\item Собрать очередь всех дуг.
\item Пока очередь не пуста:
\begin{enumerate}
\item извлечь дугу $x\to y$;
\item пересмотреть домен $D_x$;
\item удалить все значения, не имеющие поддержки в $D_y$;
\item если домен изменился, вернуть в очередь дуги, ведущие в $x$.
\end{enumerate}
\item Если какой-то домен опустел, задача несовместна.
\end{enumerate}

\subsubsection{Минимальный пример}
Пусть
\[
D_x=\{1,2\},\qquad D_y=\{2\},
\]
ограничение:
\[
x<y.
\]
Тогда:
\begin{itemize}
\item для $x=1$ существует поддержка $y=2$;
\item для $x=2$ поддержки нет, так как $2<2$ ложно.
\end{itemize}
Значит,
\[
D_x=\{1\}.
\]

\subsection{Эвристики выбора переменной и значения}
Практика решения CSP существенно улучшается эвристиками.

\subsubsection{Выбор переменной}
\begin{itemize}
\item \textbf{MRV} (minimum remaining values): выбирать переменную с наименьшим текущим доменом;
\item \textbf{degree heuristic}: среди них выбирать переменную, участвующую в наибольшем числе ограничений.
\end{itemize}

\subsubsection{Выбор значения}
\begin{itemize}
\item \textbf{Least constraining value}: сначала пробовать значение, которое исключает меньше всего значений у соседей.
\end{itemize}

\subsection{Локальный поиск для CSP}
Вместо систематического перебора можно использовать локальный поиск.

\subsubsection{Идея}
\begin{enumerate}
\item Построить произвольное полное присваивание.
\item Пока есть нарушенные ограничения:
\begin{enumerate}
\item выбрать одну из ``плохих'' переменных;
\item изменить её значение так, чтобы уменьшить число нарушений.
\end{enumerate}
\end{enumerate}

\subsubsection{Плюсы}
\begin{itemize}
\item часто хорошо работает на больших экземплярах;
\item не требует хранения дерева поиска.
\end{itemize}

\subsubsection{Минусы}
\begin{itemize}
\item может застревать в локальных минимумах;
\item не всегда даёт доказательство отсутствия решения.
\end{itemize}

\subsection{Сведение CSP к SAT}
Очень важный общий метод --- кодирование CSP в булеву выполнимость.

\subsubsection{Идея}
Если переменная $x$ имеет домен $\{a_1,\dots,a_k\}$, вводятся булевы переменные
\[
x_{a_1},\dots,x_{a_k},
\]
где $x_{a_j}=1$ означает ``переменная $x$ получила значение $a_j$''.

Далее добавляются формулы:
\begin{itemize}
\item каждая переменная получает хотя бы одно значение;
\item не более одного значения;
\item ограничения кодируются запретами несовместимых комбинаций.
\end{itemize}

\subsubsection{Преимущество}
После такого сведения можно применять SAT-решатели.

\subsection{Сведение CSP к 0/1-целочисленному программированию}
По аналогии можно кодировать допустимые значения бинарными переменными и записывать ограничения линейными неравенствами. Тогда задача решается как 0/1-IP.

\section{Эволюционные алгоритмы}

\subsection{Общая идея}
Эволюционные алгоритмы --- это семейство методов оптимизации, вдохновлённых биологической эволюцией. Вместо одного решения рассматривается \textbf{популяция} решений, которая постепенно изменяется и улучшается за счёт:
\begin{itemize}
\item отбора;
\item скрещивания;
\item мутаций.
\end{itemize}

Обычно они применяются там, где:
\begin{itemize}
\item пространство решений очень велико;
\item точный поиск слишком дорог;
\item известна функция качества решения, но задача сложна для точной оптимизации.
\end{itemize}

\subsection{Постановка задачи оптимизации}
Эволюционные алгоритмы обычно применяются к задачам оптимизации следующего вида:
\[
\text{оптимизировать } F(x)
\]
при ограничениях
\[
g_i(x)\le 0,\quad i=1,\dots,m,
\]
\[
h_j(x)=0,\quad j=1,\dots,s,
\]
где:
\begin{itemize}
\item $x$ --- искомое решение;
\item $F(x)$ --- \textbf{целевая функция}, которую нужно минимизировать или максимизировать;
\item $g_i(x), h_j(x)$ --- ограничения;
\item допустимые значения переменных задают \textbf{область поиска}.
\end{itemize}

\begin{definition}
\textbf{Допустимым решением} называется такое значение $x$, которое удовлетворяет всем ограничениям задачи.
\end{definition}

\begin{definition}
\textbf{Недопустимым решением} называется такое значение $x$, которое нарушает хотя бы одно ограничение.
\end{definition}

Таким образом, задача эволюционной оптимизации состоит в том, чтобы найти допустимое решение $x$, для которого значение целевой функции $F(x)$ является наилучшим.

\subsection{Основные понятия}
\begin{definition}
\textbf{Особь} --- одно кандидатное решение задачи.
\end{definition}

\begin{definition}
\textbf{Популяция} --- набор особей.
\end{definition}

\begin{definition}
\textbf{Генотип} --- кодировка решения (например, битовая строка, перестановка, вектор чисел).
\end{definition}

\begin{definition}
\textbf{Функция приспособленности} (fitness) --- числовая оценка качества решения.
\end{definition}

На практике функция приспособленности обычно строится на основе целевой функции $F(x)$, а также может учитывать штрафы за нарушение ограничений.

\subsection{Базовая схема генетического алгоритма}
\textbf{Вход:} задача оптимизации, целевая функция, ограничения.

\textbf{Шаги:}
\begin{enumerate}
\item Сгенерировать начальную популяцию.
\item Вычислить приспособленность всех особей.
\item Пока не выполнен критерий остановки:
\begin{enumerate}
\item выбрать родителей;
\item применить оператор скрещивания;
\item применить оператор мутации;
\item вычислить приспособленность потомков;
\item сформировать новую популяцию.
\end{enumerate}
\item Вернуть лучшую найденную особь.
\end{enumerate}

\subsection{Минимальный пример}
Пусть нужно максимизировать число единиц в битовой строке длины $4$:
\[
f(x)=\text{число единиц в }x.
\]

\subsubsection*{Начальная популяция}
\[
0001,\quad 1010,\quad 1110,\quad 0011.
\]
Их значения:
\[
1,\ 2,\ 3,\ 2.
\]

\subsubsection*{Отбор}
С большей вероятностью выбираем строку $1110$.

\subsubsection*{Скрещивание}
Пусть выбраны родители
\[
1110,\qquad 1010.
\]
При одноточечном скрещивании после второй позиции получаем потомков:
\[
11|10,\ 10|10 \;\Longrightarrow\; 1110,\ 1010.
\]

\subsubsection*{Мутация}
Например,
\[
1010\to 1011.
\]
Приспособленность увеличилась с $2$ до $3$.

\subsection{Представление решений}
Выбор кодировки крайне важен.

\subsubsection{Битовая кодировка}
Удобна для булевых задач и задач выбора подмножества.

Пример: в задаче выбора объектов строка
\[
101001
\]
может означать, что выбраны объекты $1,3,6$.

\subsubsection{Целочисленный или вещественный вектор}
Применяется для параметрической оптимизации.

\subsubsection{Перестановка}
Используется для маршрутов, расписаний и задач упорядочивания, например для коммивояжёра.

\subsection{Отбор}
Цель отбора --- \textbf{получить новое поколение, которое приведёт к улучшению решения}, то есть повысит качество популяции по значению функции приспособленности.

\subsubsection{Типовые методы}
\begin{itemize}
\item пропорциональный отбор;
\item турнирный отбор;
\item ранговый отбор;
\item рулетка.
\end{itemize}

\subsubsection{Турнирный отбор}
\begin{enumerate}
\item случайно выбрать несколько особей;
\item из них взять лучшую.
\end{enumerate}

\textbf{Плюс:} простота.
\textbf{Минус:} слишком сильный отбор уменьшает разнообразие популяции.

\subsection{Скрещивание}
Скрещивание комбинирует признаки двух родителей.

\subsubsection{Одноточечное скрещивание}
Для битовых строк:
\[
x=x_1x_2\ldots x_n,\qquad y=y_1y_2\ldots y_n.
\]
Выбирается точка $k$, потомки:
\[
x_1\ldots x_k y_{k+1}\ldots y_n,
\qquad
y_1\ldots y_k x_{k+1}\ldots x_n.
\]

\subsubsection{Двухточечное скрещивание}
Меняется средний фрагмент.

\subsubsection{Специальное скрещивание для перестановок}
Если решение есть перестановка, обычное побитовое скрещивание даёт недопустимые объекты, поэтому используются специальные операторы, сохраняющие перестановочную структуру.

\subsection{Мутация}
Мутация вводит случайные изменения и препятствует преждевременному вырождению популяции.

\subsubsection{Примеры}
\begin{itemize}
\item инверсия одного бита;
\item случайное изменение координаты;
\item перестановка двух элементов маршрута.
\end{itemize}

\subsection{Формирование следующего поколения}
Варианты:
\begin{itemize}
\item полная замена популяции потомками;
\item смешанная схема ``родители + потомки'';
\item элитизм --- часть лучших особей гарантированно переносится дальше.
\end{itemize}

\subsection{Критерии остановки}
\begin{itemize}
\item достигнуто заданное число поколений (достигнут лимит шагов);
\item определены верхняя/нижняя граница значения целевой функции и число eps $(|f(x) - bound| < \epsilon)$;
\item долго нет улучшения $(|f_i - f_{i-1}| < \epsilon)$;
\item исчерпан лимит времени.
\end{itemize}

\subsection{Корректность в строгом и практическом смысле}
Для эволюционных алгоритмов обычно различают:
\begin{itemize}
\item \textbf{формальную корректность}: алгоритм всегда возвращает допустимый объект, если кодировка и операторы построены корректно;
\item \textbf{оптимальность}: обычно точный глобальный оптимум не гарантируется.
\end{itemize}

Таким образом, эволюционные алгоритмы --- это \textbf{приближённые алгоритмы}: они ориентированы на поиск хорошего, близкого к оптимальному решения за приемлемое время.

\subsection{Доказательство корректности общей схемы}
\begin{proposition}
Если:
\begin{enumerate}
\item начальная популяция состоит из допустимых решений;
\item операторы скрещивания и мутации сохраняют допустимость;
\item функция отбора выбирает только допустимые решения,
\end{enumerate}
то любая особь на любом поколении допустима.
\end{proposition}

\begin{proof}
По индукции по номеру поколения.

\textbf{База.} На нулевом поколении все особи допустимы по условию.

\textbf{Переход.} Предположим, в текущем поколении все особи допустимы. Отбор выбирает допустимых родителей. Скрещивание и мутация по условию сохраняют допустимость. Следовательно, все потомки допустимы. Значит, и следующее поколение состоит из допустимых решений.

Индукция завершена.
\end{proof}

\subsection{Преимущества эволюционных алгоритмов}
\begin{itemize}
\item не не накладывают ограничений на свойства функции;
\item позволяют учитывать ограничения;
\item легко адаптируются под конкретную задачу;
\item высокая скорость работы.
\end{itemize}

\subsection{Недостатки}
\begin{itemize}
\item не гарантируют точный глобальный оптимум;
\item чувствительны к выбору параметров;
\item западание в локальный/минимум (алгоритм имитации отжига решает проблему).
\end{itemize}

\subsection{Связь с другими методами}
Эволюционные алгоритмы часто комбинируют:
\begin{itemize}
\item с локальным поиском;
\item с жадной инициализацией;
\item с методами восстановления допустимости;
\item с SAT/CSP-моделями как генераторами стартовых решений.
\end{itemize}

\section{Краткое сравнение трёх подходов}

\begin{center}
\begin{tabular}{|>{\raggedright\arraybackslash}p{3.7cm}|>{\raggedright\arraybackslash}p{4.1cm}|>{\raggedright\arraybackslash}p{4.1cm}|>{\raggedright\arraybackslash}p{3.6cm}|}
\hline
\textbf{Подход} & \textbf{Что даёт} & \textbf{Гарантии} & \textbf{Типичные методы} \\
\hline
Приближённые алгоритмы & Быстрое решение, близкое к оптимальному & Есть оценка отклонения от оптимума & жадные методы, LP-округление, PTAS/FPTAS \\
\hline
Методы CSP & Точное решение задачи совместимости ограничений & При полном поиске --- точность и доказательство отсутствия решения & backtracking, propagation, arc consistency, SAT-редукция \\
\hline
Эволюционные алгоритмы & Хорошее эвристическое решение на сложном пространстве поиска & Обычно нет гарантии оптимальности & генетические алгоритмы, мутации, скрещивание, селекция \\
\hline
\end{tabular}
\end{center}

\section{Краткие формулировки для устного ответа}

Ниже --- сжатая версия, удобная для ответа на экзамене.

\begin{enumerate}
\item \textbf{Теория сложности} изучает разрешимые задачи по затратам времени и памяти.
\item \textbf{Класс $P$} --- задачи, решаемые детерминированно за полиномиальное время:
\[
P=\bigcup_{c\ge1}DTIME(n^c).
\]
\item \textbf{Класс $NP$} --- задачи, для которых можно за полиномиальное время проверить сертификат:
\[
x\in L\iff \exists u,\ |u|\le p(|x|):\ M(x,u)=1.
\]
\item Имеется вложение
\[
P\subseteq NP\subseteq EXP.
\]
\item \textbf{Полиномиальная сводимость по Карпу}: $L\le_pL'$, если существует полиномиально вычислимая функция $f$ такая, что
\[
x\in L\iff f(x)\in L'.
\]
\item \textbf{NP-полная задача} --- задача из $NP$, к которой полиномиально сводится любая задача из $NP$.
\item \textbf{Теорема Кука--Левина}: задача \textsc{SAT} NP-полна.
\item Идея доказательства: по входу $x$ строится булева формула, выполнимая тогда и только тогда, когда существует принимающее вычисление проверяющей машины на $x$.
\item Важные NP-полные задачи из конспекта:
\[
\textsc{SAT},\ \textsc{3SAT},\ \textsc{0/1-IP},\ \textsc{INDSET},\ \textsc{TSP}.
\]
\item Для доказательства NP-полноты используются цепочки сведений:
\[
\textsc{CSAT}\le_p \textsc{3SAT},\quad
\textsc{CSAT}\le_p \textsc{0/1-IP},\quad
\textsc{3SAT}\le_p \textsc{INDSET},\quad
\textsc{CSAT}\le_p \textsc{dHAMPATH}\le_p \textsc{HAMPATH}\le_p \textsc{HAMCYCLE}\le_p \textsc{TSP}.
\]
\item Если хотя бы одна NP-полная задача лежит в $P$, то
\[
P=NP.
\]
\end{enumerate}

\end{document}
